% !TEX program = xelatex
% !TEX encoding = UTF-8

%%%%%%%%%%%%%%%%%%%%%%%%%%%%%%%%%%%%%%%%%%%%%%%%%%%%%%
% 专业化学位论文 LaTeX 模板
% --- 宏包与全局设置文件 (preamble.tex) ---
% --- 版本: 4.0 (目录格式最终微调版) ---
%%%%%%%%%%%%%%%%%%%%%%%%%%%%%%%%%%%%%%%%%%%%%%%%%%%%%%



%======================================================================
% 1. 页面布局与边距 (Page Layout)
%======================================================================
\usepackage[
    a4paper,
    top=2.5cm,
    bottom=2.5cm,
    left=2.5cm,
    right=2.5cm,
    headheight=15pt
]{geometry}
\usepackage{changepage} % 引入 changepage 以支持目录页边距的调整

%======================================================================
% 2. 核心功能宏包 (Core Packages)
%======================================================================
\usepackage{amsmath, amssymb, amsfonts, amsthm}
\usepackage{bm}
\usepackage{graphicx}
\graphicspath{{figures/}}
\usepackage{subcaption}
\usepackage{float}
\usepackage[section]{placeins} % 控制浮动体，避免跨节堆叠
\captionsetup[table]{position=bottom} % 表题置于表格下方
\usepackage{booktabs}
\usepackage{array}
\usepackage{longtable}
\usepackage{setspace}
\usepackage{changepage} % 提供 adjustwidth 环境
\usepackage{url}
\usepackage[table]{xcolor} % 为 listings 的 \color{gray} 等颜色提供支持

%======================================================================
% 3. 字体与字号设置 (Fonts Setup)
%======================================================================
% 使用 XeLaTeX 时需要先加载 fontspec 再加载 xeCJK，以启用 \setmainfont 等命令
\usepackage{fontspec}
\usepackage{xeCJK}
\setmainfont{Times New Roman}[Ligatures=TeX]
\setCJKmainfont{SimSun}[BoldFont=SimHei, AutoFakeSlant=true]
\setsansfont{Arial}
\setCJKsansfont{SimHei}[AutoFakeSlant=true]
\setmonofont{Consolas}

% --- 全局正文字号与行距 ---
\newcommand{\BodyFontSize}{14pt}
\newcommand{\BodyBaseline}{22pt} % 约1.6倍行距
\makeatletter
\renewcommand\normalsize{\@setfontsize\normalsize{\BodyFontSize}{\BodyBaseline}}
\makeatother
\normalsize

% 轻度放宽断行，减少 Overfull \hbox 警告
\emergencystretch=2em

%======================================================================
% 4. 章节与标题格式 (Chapter & Section Styling)
%======================================================================
% 章节编号：标题使用中文数字（使用 zhnumber 的 \zhnumber）其余保持阿拉伯数字不变
\renewcommand{\thesection}{\zhnumber{\arabic{section}}}
\renewcommand{\thesubsection}{\arabic{section}.\arabic{subsection}}
\numberwithin{equation}{section}
% 确保公式编号中的“章”部分仍为阿拉伯数字（避免出现“\CJKnumber.1”样式）
\renewcommand{\theequation}{\arabic{section}.\arabic{equation}}
\providecommand{\frontmatter}{\clearpage\pagenumbering{Roman}}
\providecommand{\mainmatter}{\clearpage\pagenumbering{arabic}}

% --- 定理环境（用于命题、引理、证明等） ---
\newtheorem{theorem}{定理}
\newtheorem{lemma}{引理}
\newtheorem{proposition}{命题}
\theoremstyle{remark}
\newtheorem{remark}{注}

%======================================================================
% 5. 目录、图录、表录格式 (TOC/LOF/LOT Styling) - 【规范格式最终版】
%======================================================================
\usepackage{tocloft}
\usepackage[section,notbib]{tocbibind}

% --- 目录主标题格式 ---
\renewcommand{\cfttoctitlefont}{\hfill\bfseries\zihao{2}}
\renewcommand{\cftaftertoctitle}{\hfill}
\renewcommand{\cftloftitlefont}{\hfill\bfseries\zihao{2}}
\renewcommand{\cftafterloftitle}{\hfill}
\renewcommand{\cftlottitlefont}{\hfill\bfseries\zihao{2}}
\renewcommand{\cftafterlottitle}{\hfill}

% --- 目录内容格式 ---
% (1) "章" 标题格式
\renewcommand{\cftsecpresnum}{第}
\renewcommand{\cftsecaftersnum}{章}
\setlength{\cftsecnumwidth}{4.5em}

% (2) 引导点格式 (数值越小越密集)
\renewcommand{\cftdotsep}{1.5}
\renewcommand{\cftsecleader}{\cftdotfill{\cftdotsep}}
\renewcommand{\cftsubsecleader}{\cftdotfill{\cftdotsep}}
\renewcommand{\cftsubsubsecleader}{\cftdotfill{\cftdotsep}}

% (3) 各级条目字体
\renewcommand{\cftsecfont}{\bfseries\zihao{3}}
\renewcommand{\cftsecpagefont}{\bfseries\zihao{3}}
\renewcommand{\cftsubsecfont}{\zihao{4}}
\renewcommand{\cftsubsecpagefont}{\zihao{4}}
\renewcommand{\cftsubsubsecfont}{\zihao{-4}}
\renewcommand{\cftsubsubsecpagefont}{\zihao{-4}}

% (4) 各级条目缩进 (实现节标题与章标题的左对齐)
\setlength{\cftsubsecindent}{\cftsecnumwidth} % 节的缩进量 = 章编号的宽度
\addtolength{\cftsubsecindent}{\cftsecindent}  % 再加上章本身的缩进
\setlength{\cftsubsubsecindent}{5.5em}

% (5) 垂直间距 (调小以使目录更紧凑)
\setlength{\cftbeforesecskip}{2pt}
\setlength{\cftbeforesubsecskip}{1pt}
\setlength{\cftbeforesubsubsecskip}{1pt}

% --- 为"目录"和"摘要"条目名称添加字间距 ---
\renewcommand{\contentsname}{目\hspace{1.5em}录}
\newcommand{\AbstractTOCTitle}{摘\hspace{1.5em}要}

% --- [核心] 重定义 \tableofcontents 等命令, 使其内容块收窄 ---
\let\oldtableofcontents\tableofcontents
\renewcommand{\tableofcontents}{%
    \begin{adjustwidth}{1cm}{1cm}
    \oldtableofcontents
    \end{adjustwidth}
}
\let\oldlistoffigures\listoffigures
\renewcommand{\listoffigures}{%
    \begin{adjustwidth}{1cm}{1cm}
    \oldlistoffigures
    \end{adjustwidth}
}
\let\oldlistoftables\listoftables
\renewcommand{\listoftables}{%
    \begin{adjustwidth}{1cm}{1cm}
    \oldlistoftables
    \end{adjustwidth}
}

%======================================================================
% 6. 页眉页脚 (Header & Footer)
%======================================================================
\usepackage{fancyhdr}
\pagestyle{fancy}
\fancyhf{}
\fancyhead[L]{\small\leftmark}
\fancyfoot[C]{\thepage}
\renewcommand{\headrulewidth}{0.4pt}
\renewcommand{\footrulewidth}{0pt}
\renewcommand{\sectionmark}[1]{\markboth{#1}{}}
\renewcommand{\subsectionmark}[1]{}
\fancypagestyle{plain}{\fancyhf{}\fancyhead[L]{\small \leftmark}\cfoot{\thepage}\renewcommand{\headrulewidth}{0.4pt}}

%======================================================================
% 7. 摘要环境定义 (Abstract Environment)
%======================================================================
\renewenvironment{abstract}{%
    \clearpage
    \thispagestyle{plain}
    \begin{center}
        \bfseries\zihao{2} 摘\hspace{1em}要
    \end{center}
    % 标题下方留白略减小
    \vspace{0.5em}
    % 摘要正文首行缩进
    \setlength{\parindent}{2em}
    % 仅对摘要启用更紧凑的行距与更宽的文字区域
    \begingroup
    \normalsize
    \setlength{\baselineskip}{18pt} % 行距从全局 22pt 收紧到 18pt
    \begin{adjustwidth}{-0.5cm}{-0.5cm} % 左右各放宽 0.5cm
    % 保证摘要内每一自然段首行缩进（安全实现）：
    \setlength{\parskip}{0pt}% 段间距为0，保持标准缩进行为
    \indent % 让首段也缩进；其余段落由 \parindent=2em 自动缩进
}{%
    \end{adjustwidth}
    \endgroup
    \par\vfill\clearpage
}

%======================================================================
% 8. 图表标题 (Caption Styling)
%======================================================================
\usepackage{caption}
\captionsetup{font=small, labelfont=bf, justification=centering}

%======================================================================
% 9. 代码块格式 (Code Listing)
%======================================================================
\usepackage{listings}
\lstset{
    language=Python, basicstyle=\footnotesize\ttfamily, keywordstyle=\bfseries,
    commentstyle=\itshape\color{gray}, stringstyle=\ttfamily,
    showstringspaces=false, breaklines=true, frame=single, framerule=0.4pt,
    numbers=left, numberstyle=\tiny\color{gray}, captionpos=b,
    tabsize=4, keepspaces=true, columns=flexible,
    literate=*{∈}{{\ensuremath{\in}}}1 {⊕}{{\ensuremath{\oplus}}}1 {Δ}{{\ensuremath{\Delta}}}1 {Σ}{{\ensuremath{\Sigma}}}1 {≤}{{\ensuremath{\le}}}1 {≥}{{\ensuremath{\ge}}}1
}



%======================================================================
% 10. 超链接与交叉引用 (Hyperlinks & Cross-references)
%======================================================================
\usepackage{hyperref}
\hypersetup{
    colorlinks=true, linkcolor=black, citecolor=black, urlcolor=black,
    pdfauthor={}, pdftitle={关于农作物种植策略的优化模型}
}
\usepackage[nameinlink, noabbrev]{cleveref}
% 使用 \crefformat 正确生成 “第n章”，避免在 \csname 构造中混入中文造成解析问题
\crefformat{section}{第#2#1#3章}
\Crefformat{section}{第#2#1#3章}
\crefname{figure}{图}{~}
\Crefname{figure}{图}{~}
\crefname{table}{表}{~}
\Crefname{table}{表}{~}
\crefname{equation}{公式}{~}
\Crefname{equation}{公式}{~}

% 为算法环境添加中文名称，便于 \cref 引用
\crefname{algorithm}{算法}{算法}
\Crefname{algorithm}{算法}{算法}

%======================================================================
% 10+. 参考文献（使用 biblatex + biber）
%======================================================================
\usepackage[backend=biber,style=numeric,sorting=none]{biblatex}
\addbibresource{references.bib}

%======================================================================
% 10+. 算法伪代码 (Algorithms)
%======================================================================
\usepackage{algorithm}
\usepackage{algpseudocode}

% 绘图（用于简洁流程图）
\usepackage{tikz}
\usetikzlibrary{arrows.meta, positioning, shapes.geometric}

%======================================================================
% 11. 论文全局信息 (Thesis Info)
%======================================================================
\title{\heiti \zihao{2} \bfseries 关于农作物种植策略的优化模型}
\author{}
\date{}